% 来源：OEBPS/Text/01_22.xhtml（OPF spine 顺序；仅第一册）
\cleardoublepage
\BookEntry[贾雷德·戴蒙德《崩溃》]{贾雷德·戴蒙德《崩溃》}

\BookCover{images/书单09崩溃.jpg}{[美] 贾雷德·戴蒙德著| 叶臻、江滢译| 上海译文出版社2011年版}

\begin{BookSummaryBox}
{\centering \textbf{美国国家科学院院士贾雷德·戴蒙德\\ 探究人类社会与文明的著作}\par}
\end{BookSummaryBox}

\subsubsection*{这本书谈什么？}

在复杂的环境问题面前，为什么有些人类社会崩溃了？有些却成功存活下来？正是带着这样的问题，书中以对失败的比较案例研究，试图为当今的人类社会提供一条生存与发展之道。
\par

在“过去社会”部分，该书讲述了一个个发人深省的失败故事。其中，有对复活节岛和玛雅文明这样经典案例的创新分析，也有对维京人殖民浪潮中两种截然不同命运的精彩叙述：冰岛的维京社会克服恶劣环境，成为当今世界人均收入最高的国家之一；而格陵兰岛上的维京社会，却因为错误的决策，在与土著人的竞争和挑战面前溃不成军。
\par

在“现代社会”部分，作者分析了发生在卢旺达的种族大屠杀，并跳脱了宗族仇恨和宗教歧视的框框，将其放入一个小型社会发生全面崩溃的视角下，加以全新解读。对于生活在同一个岛屿上的两个民族：多米尼加和海地，两种不同的国家命运给第三世界国家的发展思考提供了启示。
\par

最后，书中用大量生动的案例，讲述了跨国大企业与生态环境保护之间的微妙关系，作为“实践教训”部分的总结。
\par

\subsubsection*{芒格为什么谈起这本书？}

芒格在“专业人士需要更多的跨学科技能”的演讲中谈道：
\par

\begin{quotation}
我很早就学过基本物理学和数学，我花了很多精力，掌握了硬科学的基本治学精神与方法，我用这种方法去学习各种软科学，增加我的跨学科普世智慧。因而，我的生活无意间成了一种教育实验：一个很好地掌握了自己专业的人在非常广阔的学术领域试验这种几本治学精神与方法的可行性和有效性。
\par
\end{quotation}

实践跨学科的一个典型就是贾雷德·戴蒙德。1999年，克林顿总统为他颁发国家科学奖章，以表彰他在进化生物学领域取得的突出领域，以及利用达尔文理论在生理学、生态学等不同领域中取得的优秀成果。贾雷德·戴蒙德原本是一名生理学家，后来是洛杉矶加利福尼亚大学的地理学教授，他出色地综合了许多学科的伟大模型。《崩溃：社会如何选择成败兴亡》就是这位普利策奖获得者的作品。
\par

\subsubsection*{《崩溃》作者贾雷德·戴蒙德如是说}

“就拿森林退化这一问题来说，新几内亚高地、日本、蒂科皮亚岛和汤加岛等都曾遇到过，它们采用成功的森林管理方法化险为夷；相形之下，复活节岛、波利尼西亚群岛中的芒阿雷瓦及格陵兰的维京社会却因为没能找到相应对策而走向消亡。我们是如何理解这些不同的结局呢？一个社会的应对之道取决于其政治、经济、社会制度和文化价值观，这些制度和价值观影响到社会能否解决其面对的问题（或是否尝试去解决问题）。”
\par

“悲观主义者会注意到中国存在大量危险，其一是对经济增长的重视程度高于环境保护和可持续发展；其二，大众的环境保护意识低下，这部分是因为教育投入不够多，中国的教育经费占国民生产总值的比重不及第一世界国家的一半。中国的人口总数占世界人口的20\%，但教育经费只是全世界教育经费的1\%。”
\par
